\chapter{Introducción}
\label{chap:introduccion}

\section{Contexto y motivación}
\label{sec:contexto-motivacion}

El acceso a la información jurídica es una condición necesaria para que la ciudadanía pueda conocer el marco normativo que afecta a su vida cotidiana. Las normas regulan aspectos tan diversos como las relaciones laborales, los derechos de los consumidores, los procedimientos administrativos, la protección de datos, la vivienda, la movilidad, las ayudas públicas o la actividad económica. Sin embargo, disponer de textos normativos publicados no implica necesariamente que esos textos sean fáciles de localizar, interpretar de forma preliminar o relacionar con una duda formulada en lenguaje ordinario.

En España, el \emph{Boletín Oficial del Estado} constituye una fuente pública esencial para la publicación y consulta de normas. Además de la publicación oficial de disposiciones, la Agencia Estatal Boletín Oficial del Estado ofrece servicios de legislación consolidada que integran en un único texto las modificaciones expresas que ha sufrido una norma desde su aprobación inicial. Esta consolidación facilita la consulta práctica de la normativa, porque evita que el usuario tenga que reconstruir manualmente la versión aplicable a partir de sucesivas reformas. No obstante, la propia Agencia advierte de que los textos consolidados tienen carácter meramente informativo y carecen de valor jurídico, por lo que la fuente con plena validez sigue siendo el diario oficial correspondiente \cite{boeLegislacionActualizada,boeApiConsolidada}.

La disponibilidad digital de estos recursos permite plantear soluciones tecnológicas orientadas a mejorar la búsqueda y exploración de información jurídica. En el ámbito europeo, iniciativas como el Portal Europeo de e-Justicia persiguen precisamente facilitar el acceso a información sobre sistemas judiciales y mejorar el acceso a la justicia en varios idiomas \cite{ejusticePortal}. De forma más amplia, organismos internacionales han señalado que los sistemas digitales y las técnicas de inteligencia artificial pueden contribuir a mejorar el acceso a información y servicios jurídicos, siempre que se diseñen con criterios de transparencia, control humano, seguridad y respeto a los derechos fundamentales \cite{oecdAccessJusticeAI,cepejEthicalCharter}.

En este contexto, los modelos de lenguaje de gran tamaño han mostrado una capacidad notable para procesar instrucciones en lenguaje natural y generar respuestas coherentes. Trabajos como el de Brown et al. \cite{brown2020language} evidenciaron que el aumento de escala de estos modelos puede mejorar su rendimiento en tareas variadas mediante aprendizaje con pocos ejemplos. Sin embargo, esta capacidad no elimina limitaciones críticas en dominios sensibles. Los modelos generativos pueden producir respuestas fluidas pero incorrectas, desactualizadas o no verificables. Estas limitaciones han sido ampliamente discutidas en la literatura sobre riesgos de los modelos de lenguaje \cite{bender2021stochastic} y adquieren especial relevancia en el ámbito jurídico, donde una afirmación aparentemente convincente puede inducir a error si no está respaldada por una fuente normativa concreta.

Una línea de trabajo especialmente relevante para este problema es la generación aumentada mediante recuperación, conocida como \emph{Retrieval-Augmented Generation} (RAG). Este enfoque combina un modelo generativo con un mecanismo de recuperación de documentos o fragmentos relevantes, de modo que la respuesta se construya a partir de información externa recuperada en el momento de la consulta. Lewis et al. \cite{lewis2020rag} mostraron que este tipo de arquitectura permite combinar memoria paramétrica y memoria no paramétrica, mejorando la especificidad y factualidad de las respuestas en tareas intensivas en conocimiento. Para un sistema aplicado a normativa, esta idea resulta especialmente atractiva: en lugar de pedir al modelo que responda únicamente desde sus parámetros internos, se le proporciona un contexto seleccionado a partir de documentos jurídicos indexados.

El presente Trabajo Fin de Grado se sitúa en esta intersección entre acceso a información jurídica, recuperación de información y modelos generativos. Su motivación principal es diseñar e implementar una aplicación demostrable que permita formular consultas en lenguaje natural sobre un corpus de normativa consolidada procedente del BOE, recuperar evidencias documentales relevantes y generar respuestas trazables. La finalidad no es sustituir la lectura de la fuente oficial ni ofrecer asesoramiento jurídico profesional, sino facilitar una primera aproximación informativa que ayude al usuario a localizar los fragmentos normativos relacionados con su consulta.

\section{Dificultad de consultar normativa extensa y cambiante}
\label{sec:dificultad-consulta-normativa}

La consulta de normativa presenta dificultades que no se resuelven únicamente con la digitalización de los documentos. En primer lugar, los textos jurídicos suelen ser extensos y estar organizados mediante estructuras jerárquicas complejas: títulos, capítulos, secciones, artículos, disposiciones adicionales, transitorias, derogatorias y finales. Una pregunta aparentemente sencilla puede depender de varios artículos o de disposiciones situadas en lugares distintos de una misma norma. Además, el significado práctico de un precepto puede depender de definiciones, excepciones o remisiones internas que no siempre son evidentes para un lector no especializado.

En segundo lugar, existe una distancia entre el lenguaje jurídico y el lenguaje con el que la ciudadanía formula sus dudas. Una persona puede preguntar por una ``ayuda'', una ``multa'', una ``baja'', una ``reclamación'' o un ``plazo'', mientras que el texto normativo puede utilizar términos como prestación, sanción, incapacidad temporal, recurso administrativo o cómputo de días hábiles. Esta diferencia terminológica dificulta la recuperación basada exclusivamente en coincidencias literales y justifica el uso de técnicas de recuperación semántica capaces de relacionar consultas y fragmentos aunque no compartan exactamente las mismas palabras.

En tercer lugar, la normativa es dinámica. Las normas pueden modificarse, derogarse, corregirse o quedar afectadas por disposiciones posteriores. La legislación consolidada del BOE ayuda a consultar una versión integrada del texto, pero no elimina la necesidad de registrar qué versión se ha utilizado, cuándo se ha capturado y cuál era su estado de actualización. La API de legislación consolidada permite recuperar metadatos, textos completos, versiones y bloques documentales, lo que resulta adecuado para construir un corpus reproducible; aun así, el sistema debe conservar información suficiente para que los resultados puedan interpretarse en relación con una versión documental concreta \cite{boeApiConsolidada}.

Estas características convierten la consulta normativa en una tarea de recuperación de información especializada. No basta con localizar documentos que contengan algunas palabras de la pregunta: es necesario recuperar fragmentos jurídicamente pertinentes, mantener su contexto y permitir que el usuario inspeccione la fuente de la respuesta. La literatura reciente sobre inteligencia artificial jurídica confirma que los modelos de lenguaje aplicados al Derecho requieren evaluación específica, porque las capacidades generales de un modelo no garantizan fiabilidad en tareas jurídicas. LegalBench, por ejemplo, se propuso como un banco de pruebas colaborativo para medir distintos tipos de razonamiento jurídico en modelos de lenguaje \cite{guha2023legalbench}. Asimismo, estudios sobre alucinaciones legales muestran que los modelos pueden generar afirmaciones incompatibles con hechos jurídicos o referencias incorrectas, lo que refuerza la necesidad de mecanismos de verificación y evaluación \cite{dahl2024legalhallucinations}.

La arquitectura RAG reduce parte de este riesgo al fundamentar la generación en documentos recuperados, pero no lo elimina. Un sistema de este tipo puede fallar si recupera fragmentos irrelevantes, si omite información necesaria, si el contexto entregado al modelo es insuficiente o si la respuesta generada no respeta fielmente las evidencias disponibles. Evaluaciones recientes de herramientas jurídicas basadas en RAG han señalado precisamente que la recuperación documental, la generación y la verificación de citas deben analizarse conjuntamente \cite{magesh2025hallucinationfree,kim2025legalragfail}. Por este motivo, en este TFG la trazabilidad no se considera un complemento visual, sino una propiedad central del sistema: cada respuesta debe estar vinculada a evidencias recuperadas que permitan revisar de dónde procede la información.

\section{Planteamiento del problema}
\label{sec:planteamiento-problema}

El problema abordado en este Trabajo Fin de Grado puede formularse de la siguiente manera: diseñar, implementar y evaluar una aplicación basada en una arquitectura RAG capaz de responder consultas formuladas en lenguaje natural sobre un corpus ampliable de normativa consolidada procedente del BOE, recuperando evidencias relevantes y generando respuestas trazables que faciliten su revisión por parte del usuario.

Esta formulación combina un problema de acceso a la información con un problema técnico de recuperación, generación y evaluación. Desde el punto de vista del usuario, la necesidad consiste en poder expresar una duda de forma natural y recibir una respuesta comprensible, acompañada de referencias a los fragmentos normativos utilizados. Desde el punto de vista del sistema, la dificultad consiste en transformar documentos jurídicos extensos en una representación indexable, recuperar información pertinente para cada consulta, construir un contexto controlado y generar una respuesta que no exceda lo respaldado por las evidencias disponibles.

El enfoque elegido se basa en RAG porque permite separar dos funciones complementarias. Por un lado, el recuperador selecciona fragmentos del corpus que pueden contener la información necesaria. Por otro, el modelo generativo utiliza esos fragmentos para redactar una respuesta en lenguaje natural. Esta separación es importante en un dominio como el jurídico: permite actualizar o ampliar el corpus sin depender únicamente del conocimiento interno del modelo y proporciona una base para inspeccionar las fuentes utilizadas. Además, facilita la evaluación modular del sistema, ya que pueden analizarse por separado la calidad de la recuperación, la adecuación del contexto construido, la fidelidad de la respuesta y la corrección de las citas.

El proyecto se orienta principalmente a la ciudadanía general, aunque también puede resultar útil para profesionales, personal de administraciones u otros usuarios que necesiten explorar documentación normativa. No obstante, el sistema se diseña como una herramienta informativa y de apoyo a la consulta documental. No pretende emitir dictámenes jurídicos, resolver casos particulares con efectos legales ni sustituir el criterio de profesionales del Derecho. Esta delimitación es necesaria tanto por la naturaleza informativa de la legislación consolidada como por las limitaciones propias de los sistemas automáticos de generación de texto.

\section{Alcance y limitaciones}
\label{sec:alcance-limitaciones}

El alcance del trabajo comprende la construcción de una aplicación demostrable y extensible para consultar normativa mediante lenguaje natural. El sistema debe permitir incorporar un corpus seleccionado de legislación consolidada, procesar los documentos, preservar metadatos relevantes, fragmentar el contenido, generar representaciones vectoriales, recuperar evidencias ante una consulta, construir un contexto de entrada para el modelo generativo y producir una respuesta acompañada de referencias verificables. La memoria documentará tanto las decisiones de diseño como la evaluación experimental del sistema.

El corpus inicial del proyecto se concibe como una colección ampliable de normas consolidadas procedentes del BOE. En la introducción no se fija una cifra definitiva de documentos, porque el tamaño exacto del corpus puede variar durante el desarrollo. La versión utilizada para los experimentos se documentará de forma explícita en los capítulos de implementación y validación, indicando la fecha de captura, los criterios de selección, los identificadores de las normas, el estado de consolidación y las estadísticas resultantes del procesamiento. Esta decisión permite mantener la introducción estable y, al mismo tiempo, asegurar la reproducibilidad de los resultados experimentales.

Quedan fuera del alcance de este TFG la construcción de un sistema jurídico de producción, la cobertura completa del ordenamiento jurídico español, la emisión de asesoramiento legal personalizado, la sustitución de fuentes oficiales y la validación jurídica exhaustiva por especialistas del Derecho. También se consideran trabajo futuro la incorporación automática y periódica de nuevas normas, la consulta sobre jurisprudencia o doctrina, la evaluación con usuarios reales a gran escala y el despliegue con requisitos completos de seguridad, monitorización y mantenimiento.

El sistema presenta además limitaciones técnicas y metodológicas. La calidad de las respuestas dependerá de la cobertura y calidad del corpus indexado, de la estrategia de fragmentación, del modelo de embeddings, del recuperador, del presupuesto de contexto disponible y del comportamiento del modelo generativo. Incluso con recuperación documental, pueden producirse errores por evidencias incompletas, fragmentos mal seleccionados, ambigüedad de la consulta o generación no fiel al contexto. Por ello, el trabajo incorporará una evaluación específica de recuperación, generación, citas y análisis de errores, evitando presentar la aplicación como una fuente autónoma de verdad jurídica.

\section{Estructura de la memoria}
\label{sec:estructura-memoria}

La memoria se organiza en varios capítulos que siguen la relación entre problema, objetivos, diseño, implementación, evaluación y conclusiones. Tras esta introducción, el Capítulo~\ref{chap:objetivos} define el objetivo general del trabajo, los objetivos específicos, los requisitos funcionales y no funcionales, las restricciones del proyecto y la planificación temporal. El Capítulo~\ref{chap:estado-arte} presenta el estado del arte necesario para contextualizar el sistema, incluyendo recuperación de información, modelos de lenguaje, arquitecturas RAG, embeddings, fragmentación documental, evaluación de sistemas RAG y particularidades del dominio jurídico.

El Capítulo~\ref{chap:implementacion} describe el diseño y la implementación del sistema desarrollado. En él se detallan las fuentes documentales, el procesamiento del corpus, la generación de fragmentos y metadatos, la indexación, los mecanismos de recuperación, la construcción del contexto, la generación de respuestas y las decisiones adoptadas para preservar trazabilidad y reproducibilidad. El Capítulo~\ref{chap:experimentos} expone los experimentos realizados para validar la solución, las métricas empleadas, los resultados obtenidos y el análisis de errores.

El Capítulo~\ref{chap:impacto} analiza el impacto social, económico y medioambiental del trabajo, así como las responsabilidades éticas y profesionales asociadas a un sistema de inteligencia artificial aplicado a información jurídica. Finalmente, el Capítulo~\ref{chap:conclusiones} resume el grado de cumplimiento de los objetivos, las competencias aplicadas y adquiridas durante el desarrollo, las principales limitaciones del sistema y las líneas de trabajo futuro. La memoria se completa con la bibliografía y los anexos necesarios para documentar el repositorio, la instalación, la ejecución y otros materiales complementarios.
